\section{Invertibility and Isomorphisms}

\subsection{Invertible Linear Maps}
  \setcounter{thm}{58}
  %\ 3.59
  \begin{mydef} [invertible, inverse]
    $S\in \linmap(W,V)$ with $ST=I_V \myand TS = I_W$ is called the \qt{inverse} of the \qt{invertible} linear map $T \in \linmap(V,W)$.
  \end{mydef}

  % 3.60
  \begin{thm} [inverse is unique]
    An invertible  linear map has a unique inverse.
  \end{thm}
  \begin{prf}
    Suppose $T\in \lvw$ is invertible and $S_1$ and $S_2$ are inverses of $T$. Hence,
    \begin{equation}
      S_1 = S_1 I = S_1 (T S_2) = (S_1 T) S_2 = I S_2 = S_2.
    \end{equation}

    Therefore, $S_1 = S_2$.
  \end{prf}

  % 3.61
  \begin{mydef} [notation $T^{-1}$]
    Since the inverse is unique, the inverse is denoted by $T^{-1}\in \linmap(W,V)$ for $T\in \linmap(V,W).$ We have
    \begin{equation}
      T^{-1}T=I_V$ and $T T^{-1} = I_W.
    \end{equation}
  \end{mydef}

  %\setcounter{thm}{62}
  %\textbf{3.63}
  \begin{thm} [invertibility, injectivity and surjectivity]
    A linear map is invertible $\iff$ it is injective and surjective.
  \end{thm}

  % 3.65
  \setcounter{thm}{64}
  \begin{thm} [injectivity is equivalent to surjectivity]
    \label{thm: injectivity is equivalent to surjectivity}
    For $T\in \linmap(V,W)$, if $V$ and $W$ are finite-dimensional, injectivity is equivalent to surjectivity:
    \begin{equation}
      T$ is invertible $\iff$ $T$ is injective $\iff$ $T$ is surjective.$
    \end{equation}

    Equivalently:
    \begin{equation}
      T$ is not invertible $\iff$ $T$ is not injective $\iff$ $T$ is not surjective.$
    \end{equation}
  \end{thm}
  \begin{prf}
    The rank-nullity theorem (\ref{rank-nullity-theorem}) or fundamental theorem of linear maps states that
    \begin{equation}
      \label{rank-nullity-nested-equation}
      \dim V = \dim \mynull T + \dim \myrange T
    \end{equation}

    If $T$ is injective, $\dim \mynull T = 0$ and therefore
    \begin{equation}
      \dim \myrange T = \dim V = \dim W,
    \end{equation}
    which means $T$ is surjective ($\myrange T = \dim W$). This is because every linearly independent list of vectors of of length $\dim W$ is a basis by
    \ref{thm: linearly independent list of the right length is a basis}. And thats how dimension is defined. (TODO)

    If $T$ is surjective, we have $\dim \myrange T = \dim W$ to begin with and therefore \eqref{rank-nullity-nested-equation} becomes
    \begin{equation}
      \dim \mynull T = \dim V - \dim \myrange T = \dim V - \dim W = 0.
    \end{equation}

    which makes $T$ also injective.
    Since injectivity and surjectivy together imply invertibility, this ends the proof.
  \end{prf}


  % 3.68
  \setcounter{thm}{67}
  \begin{thm} [$ST = I \iff TS=I$ on vector spaces of the same dimension]
    \label{thm: ST = I iff TS=I - on vector spaces of the same dimension}
    Suppoe $V$ and $W$ are finite-dimensinal vector spaces of the same dimension. Let $S \in \linmap (W,V)$, and $T \in \linmap(V,W)$. Then $ST = I$ $\iff$ $TS = I$.
  \end{thm}
  \begin{prf}
     \Rightarrowdirection Let $S \in \linmap (W,V)$, and $T \in \linmap(V,W)$ such that $ST = I$. If $v \in V$ and $Tv = 0$, then
     \begin{equation}
       v = Iv = (ST)v = S(Tv) = S(0) = 0.
     \end{equation}

     Thus, $\mynull T = \{0\}$ and therefore, $T$ is injective by \ref{thm: injectivity iff null space equals zero-set}. Because $V$ and $W$ have the same dimension, $T$ is invertible by \ref{thm: injectivity is equivalent to surjectivity}. Now, multiply both sides of the equation $ST = I$ by $T^{-1}$ to conclude:
     \begin{equation}
       S = T^{-1}.
     \end{equation}

     Hence, $TS = TT^{-1} = I$, as desired.

     \Leftarrowdirection Right above, we have shown that $ST = I$ implies that $TS=I$. To prove the implication in the other direction, simply reverse the roles of $S$ and $T$ and $V$ and $W$, showing that $TS=I$ implies that $ST=I$.
  \end{prf}

  \subsection{Isomorphic Vector Spaces}

  %\textbf{3.69}
  \begin{mydef}
    An isomorphism is an invertible linear map. Two vector spaces are called isomorphic if there is a isomorphism from one vector space to the other. $T:V\to W$
  \end{mydef}

  %\textbf{3.70}
  \begin{thm}[dimension shows whether vector spaces are isomorphic]
    \label{thm: dimension shows whether vector spaces are isomorphic}
    Two \fd vector spaces $U$, $V$ over $\mathbb{F}$ are isomorphic $\iff$ they have the same dimension (i.e., $\dim U = \dim V$).
  \end{thm}
  \begin{prf}
    exercise

    \Rightarrowdirection rank-nullity theorem

    \Leftarrowdirection surjectivity and injectivity.
  \end{prf}


  %\textbf{3.71}
  \begin{thm} [$\linmap (V,W)$ and $\myF^{m, n}$ are isomorphic]
    \label{thm: the space of linear maps and the space of all matrices are isomorphic}
    Suppose $\oneTillN{v}$ is a basis of $V$ and $\oneTillM{w}$ is a basis of $W$. Then $\mathcal{M}$ is an isomorphism between $\lvw$ and $\mathbb{F}^{m,n}$.
  \end{thm}
  \begin{prf}
    exercise (inectivity and surjectivity)
  \end{prf}

  %\textbf{3.72}
  \begin{thm} [$\dim \lvw = (\dim V)\cdot(\dim W)$]
    \label{thm: dimension of L(V,W) = dim L (V) * dim W (W)}
    Suppose $V,W$ are \fd. Then $\lvw$ is finite-dimensional and
    \begin{equation}
      \dim \lvw = (\dim V)\cdot(\dim W).
    \end{equation}
  \end{thm}
  \begin{prf}
    %The desired result follows from \ref{thm: the space of linear maps and the space of all matrices are isomorphic}, \ref{thm: dimension shows whether vector spaces are isomorphic}, and \ref{thm: the dimension of the vector space of all m by n matrices is mn}.

    Suppose $v_1, \ddd, v_n$ is a basis of $V$ and $w_1, \ddd, w_m$ is a basis of $W$.
    By \ref{thm: the space of linear maps and the space of all matrices are isomorphic}, there is an an isomorphism between $\lin{V}{W}$ and $\myF^{m,n}$, and thus, they have the same dimension. Now, \ref{thm: the dimension of the vector space of all m by n matrices is mn} tells us that $\dim \myF^{m,n} = m  n$, therefore
    \begin{equation}
      \dim \lin{V}{W} = \dim \myF^{m,n} = m\cdot n = (\dim V)\cdot(\dim W),
    \end{equation}

    as desired.
  \end{prf}

  \subsection{Linear Maps Thought of as Matrix Multiplication}

  \begin{mydef} [matrix of a vector, $\mmatrix(u)$]
    Let $u \in V$ and $\onetillm{v}$ be a basis of $V$ such that
  \begin{equation}
    u=b_1v_1+\dots+b_mv_m. $ We define $
        \mathcal{M}(u) :=
      \left (
      \begin{matrix}
        b_1 \\ \vdots \\ b_m
      \end{matrix}
      \right ).
  \end{equation}

    The vector depends on the basis, but it is not included in the notation.
%    Recall:
%\begin{equation}
%    \mathcal{M} (T) :\equiv
%  \begin{blockarray}{cccc}
%    & v_1 \quad \cdots & v_k & \cdots \quad v_n \\
%    \begin{block}{c(ccc)}
%      w_1    & & A_{1,k} & \\
%      \vdots & & \vdots & \\
%      w_m    & & A_{m,k} & \\
%    \end{block}
%  \end{blockarray}
%\end{equation}

  \end{mydef}

  \setcounter{thm}{74}
  %\textbf{3.75}
  \begin{thm} [$\mmatrix(T)_{\mathsmaller{\bullet}, k} = \mmatrix(T v_k)$]
    Suppose $T \in \linmap(V,W)$ and $v_1, \ldots, v_n$ is a basis of $V$ and $w_1, \ldots, w_m$ is a basis of $W$. Let $1\leq k \leq n$. Then the $k^{\text{th}}$ column of $\mmatrix(T)$, which is denoted by $\mmatrix(T)_{\mathsmaller{\bullet}, k}$, equals $\mmatrix(Tv_k)$.
    \begin{equation}
      \mmatrix(T)_{\mathsmaller{\bullet}, k}=\mmatrix(Tv_k) \quad \forall k \in \{1, \ddd, n\}.
    \end{equation}
%    $\mmatrix_{\mathsmaller{\bullet}, k} = M(T v_k)$. The $k^{\text{th}}$ column of $\mmatrix_{\mathsmaller{\bullet}, k}$ equals $(A_{1,k}, \dots, A_{m,k})^\top$
  \end{thm}

  %3.76
  \begin{thm} [linear maps act like matrix multiplication]
    Suppose $T\in \lvw$ and $v\in V$. Suppose $\onetilln{v}$ is a basis of $V$ and $\onetillm{w}$ is a basis of $W$. Then
    \begin{equation}
      \mathcal{M} (Tv)=\mathcal{M}(T) \cdot \mathcal{M}(v)
    \end{equation}

    Or using different notation:
    \begin{equation}
      \mathcal{M} (Tv, (\onetilln{v}), (\onetillm{w}))=\mathcal{M}(T, (\onetilln{v}), (\onetillm{w})) \cdot \mathcal{M}(v)
    \end{equation}

  \end{thm}

  % 3.78
  \setcounter{thm}{77}
  \begin{thm} [dimension of $\myrange T$ and column rank of $\mmatrix(T)$]
    Suppose $\dim V < \infty$ and $\dim W < \infty$. Then for $T \in \linmap(V,W):$
    \begin{equation}
      \dim \myrange T = $ column rank of $ \mmatrix (T)
    \end{equation}
  \end{thm}

  \subsection{Change of basis}


  % 3.80
  \setcounter{thm}{79}
  \begin{mydef}[invertible, inverse, $A^{-1}$]
    A square matrix $A$ is called invertible, if there is some square matrix B of the same size such that
    \begin{equation}
    	AB=BA=I.
    \end{equation}

    We call $B$ the inverse of $A$ denoted by $A^{-1} :\equiv B$. Rules:
    \begin{itemize}
      \item $(A^{-1})^{-1}=A$
      \item $(AC)^{-1} = C^{-1}A^{-1}$ (Because $(AC)(C^{-1}A^{-1})=I$ and $(C^{-1}A^{-1})(AC)=I$)
    \end{itemize}
  \end{mydef}

  \mce{81}
  \begin{thm} [matrix of product of linear maps]
    \label{thm: matrix of product of linear maps2}
    Suppose $T\in \mathcal{L}(U,V)$ and $S\in \lvw$. If $\onetillm{u}$ is a basis of $U$, $\onetilln{v}$ is a basis of $V$ and $\onetill{w}{p}$ is a basis of $W$. Note that $\dim U = m$, $\dim V = n$, $\dim W = p$. Then we have:
    \begin{equation}
    \begin{aligned}
      &\mmatrix(ST, (\onetillm{u}), (\onetill{w}{p})) = \\
      &\qquad \mmatrix(S, (\onetilln{v}), (\onetill{w}{p} ))
      \cdot
      \mmatrix(T, (\onetillm{u}), (\onetilln{v}   ))
    \end{aligned}
    \end{equation}

    Or using different notation:
    \begin{equation}
      \mathcal{M}(ST) = \mathcal{M}(S) \cdot \mathcal{M}(T)
    \end{equation}
  \end{thm}
  \begin{prf}
    This is just another way of writing \ref{thm: matrix of product of linear maps1}.
  \end{prf}

  \mce{82}
  \begin{thm} [matrix of identity operator with respect to two bases]
    \label{thm: matrix of identity operator with respect to two bases}
    Suppose that $u_1, \ddd, u_m$ and $v_1, \ddd, v_m$ are bases of $V$. Then the matrices
    \begin{equation}
      \m{I, (u_1, \ddd, u_m), (v_1, \ddd, v_m)} \qquad \text{and} \qquad \m{I, (v_1, \ddd, v_m), (u_1, \ddd,u_m)}
    \end{equation}

    are invertible, and each is the inverse of the other.
  \end{thm}
  \begin{prf}
    Carefull application of \ref{thm: matrix of product of linear maps2} yields
    \begin{equation}
      \m{I_V} = I_n = \m{I_V I_V} = \m{I_V, (v_1, \ddd, v_m), (u_1, \ddd, u_m)} \cdot \m{I_V, (u_1, \ddd, u_m), (v_1, \ddd, v_m)}.
    \end{equation}

    Now, interchange the roles of the $u$'s and the $v$'s, getting
    \begin{equation}
      \m{I_V} = I_n = \m{I_V I_V} = \m{I_V, (u_1, \ddd, u_m), (v_1, \ddd, v_m)} \cdot \m{I_V, (v_1, \ddd, v_m), (u_1, \ddd, u_m)}.
    \end{equation}

    So, the product of both matrices equal to $I_n$, which is the identity in $\myF^{n,n}$.
  \end{prf}

  % 3.84
  \setcounter{thm}{83}
  \begin{thm}[change-of-basis formula:]
    Let $T \in \op(V)$. Suppose $u_1, \ddd, u_m$ and $v_1, \ddd, v_m$ are bases of $V$ and . Let
    \begin{equation}
      \begin{aligned}
        A & := \m {T, (u_1, \ddd, u_m)}, \\
        B & := \m {T, (v_1, \ddd, v_m)}, \text{ and }\\
        C & := \m {I, (u_1, \ddd, u_m), (v_1, \ddd, v_m)}. \\
      \end{aligned}
    \end{equation}

    Then
    \begin{equation}
      A = C^{-1} B C.
    \end{equation}
  \end{thm}

  % 3.86
  \setcounter{thm}{85}
  \begin{thm}
    If $v_1, \ddd, v_m$ is a basis $V$ and $T \in \op{V}$ is invertible, then
    \begin{equation}
      \mmatrix(T^{-1}) = (\mmatrix(T))^{-1},
    \end{equation}

    where both matrices are with respect to the basis $v_1, \ddd, v_m$, that is
    \begin{equation}
      \mmatrix(T^{-1}, (\onetillm{v})) =\mmatrix(T, (\onetillm{v}))^{-1}.
    \end{equation}
  \end{thm}